\subsection*{A source-geometric short-tail estimate for both A.11 supports}
The two azimuthal corrections in source (A.11) are supported in intervals
of width $3/10$ centered three and one logarithmic radial units before
the end of the $-\lambda$ plateau. Write
\[
D(y,\eta)=\frac{1}{2XE(y,\eta)^2}\int_X^\infty E(x,\eta)^2\,dx.
\]
\begin{proposition}[Remaining exterior tail]\label{prop:a11-short-tail-new}
Assume the source's unedited plateau $\ell_b=-\lambda$, followed by
a one-unit transition with $\ell_b\le0$, a $-1$ hold of length
$4\log(1/h)$, and a remaining tail with $\ell_b\le-3h/4$.
For every point of either support,
\[
D_b(y)\le\frac{93}{40}+\frac{h^7}{3}.
\]
If $E_c=E_b(1+u)$ and $|u|\le\delta<1$ at every later radius,
then
\[
D_c\le\left(\frac{1+\delta}{1-\delta}\right)^2
 \left(\frac{93}{40}+\frac{h^7}{3}\right).
\]
For $0<\lambda\le1/20$, $h\le\lambda/100$ and
$\delta=40960\lambda^{28}$ this is strictly less than $233/100$.
\end{proposition}
\begin{proof}
Use $D_b=\frac12\int_0^\infty
\exp(2\int_0^v\ell_b(y+t)dt)\,dv$.
The remainder of the plateau is at most $63/20$ units and
contributes at most $63/40$; the next transition contributes
$1/2$. The $-1$ hold contributes at most $1/4$ and reduces
the energy factor by $h^8$. Its remaining tail contributes
at most $h^8/(3h)=h^7/3$.
The corrected-to-base kernel ratio lies in
$[((1-\delta)/(1+\delta))^2,
((1+\delta)/(1-\delta))^2]$.
Finally $\delta<10^{-31}$ and the explicit rational bound
$(1+8/10000)(93/40+1/1000)<233/100$ completes the proof.
\end{proof}

\begin{theorem}[Improved conditional A.11 axial stress and cone]
\label{thm:a11-short-tail-new}
Assume the genuine source moment and pressure gauges
$S(\infty)=0$, $\Pi(\infty)=0$, together with
$|c_i|\le192\lambda^{28}$,
$|\partial_\eta c_i|\le12\lambda^{28}$,
the source bump sup bounds, nonpositive corrected future slopes,
and $E_c\le\lambda/4$ throughout the two supports.
For $0<\lambda\le1/20$ and $0<h\le\lambda/100$,
\[
\frac{|N_s|}{E_c}<\frac{51}{200}\lambda<\frac{\lambda}{2},
\qquad \frac{\lambda}{2}-\frac{|N_s|}{E_c}>
 \frac{49}{200}\lambda.
\]
If additionally $Q_s\ge\lambda/2$,
$-3\lambda/2\le\ell_c\le-\lambda/2$ and $X\ge9/\lambda$,
then with $a=2-2\ell_c$,
$p_1=XQ_s/(1-2h\eta^2)$ and $w=N_s/(E_cQ_s)$,
\[
2(1-a/p_1)^2-(a-2)w^2\ge
\frac{92197}{200000}>\frac9{20}.
\]
\end{theorem}
\begin{proof}
The exact source stress-tail identity (4.16) gives
\[
\frac{|N_s|}{E_c}\le E_c\{
4hD_c+4AP_c+2\varepsilon(D_c+P_c)\},
\quad A=\tfrac12+h,\quad
\varepsilon=\frac{2560\lambda^{28}}{1-\delta},
\]
where $P_c=-\Pi/E_c^2$. The future slope ensures
$0\le P_c\le1/2$. The proposition gives $D_c<233/100$,
so $4hD_c<1/200$, $4AP_c\le1001/1000$ and
$2\varepsilon(D_c+P_c)<1/200$.
The bracket is smaller than $51/50$, yielding the first
strict margin. For the cone,
$|w|\le51/100$, $a\le43/20$, and $p_1\ge9/2>2a$.
Thus its left side is at least
$1/2-(3/20)(51/100)^2=92197/200000$.
\end{proof}
These are conditional source-stage inequalities, not a certified
numerical realization of an entire Appendix-A parameter tuple.
They do not check the other stress-cone stages, Appendix B's
contraction, unforced startup equations, or a zero-force physical
Navier--Stokes residual. The originating OpenAI theorem includes
smooth external forcing.
